\section{Introduction}
\label{sec:introduction}

\begin{wrapfigure}{r}{0.42\textwidth}
\vspace{-1.1em}
\centering
\includegraphics[width=\linewidth,trim=0 72pt 0 0,clip]{teaser_revised.pdf}
\vspace{-0.45cm}
\caption{\textbf{A descriptive probe.} Three forecasts after two heads from
an unfamiliar coin illustrate why two observations do not identify the operative
relationship
(Appendix~\ref{app:carnival-coin-probe}).}
\vspace{-0.45cm}
\label{fig:world-knowledge-evidence}
\end{wrapfigure}

A forecast is a compact theory of how a world works. News of a
policy proposal, an endorsement, or an economic shock matters only through a
relationship: which actors can respond, which institutions transmit the effect,
and whether the local setting amplifies or dampens it. The same observation
therefore supports different predictions in different contexts. National
political news carries more predictive weight in a media environment organized
around national politics than where local issues dominate. Forecasting requires selecting the
relationship that applies before extrapolating from the evidence.
Figure~\ref{fig:world-knowledge-evidence} illustrates the identification
problem: identical observations support different forecasts under different
assumptions about the generator.

LMs now approach crowd accuracy, a practical performance floor; ensembles
match it on selected question sets \citep{halawi2024approaching,schoenegger2024wisdom}. This makes the
process behind a probability increasingly important. High average scores can
coexist with reactions to the wrong cue or the transfer of a relationship into a
setting where it no longer applies.

But accuracy alone does not resolve this problem. The same probability can come
from a base rate, a retrieved precedent, a market price, a lexical shortcut, or a
relationship inferred from the current context. Generated rationales do not
resolve the ambiguity because explanations need not track model behavior
\citep{jacovi2020towards,turpin2023language}. We therefore ask a process question:
\emph{when an LM forecasts an outcome, does it behave as if it selects,
revises, and transfers predictive structure?} The problem is central to
\textbf{judgmental forecasting}, which combines base rates, sparse observations,
qualitative evidence, and domain knowledge to estimate future events
\citep{mellers2014psychological,tetlock2016superforecasting}.

We call this behavioral process \emph{inductive reasoning}: selecting a predictive
relationship from context and sparse examples, revising that selection as local
evidence accumulates, and applying it to a new case
\citep{tenenbaum2006theory,griffiths2010probabilistic,lake2017building}. We use
\emph{world structure} for the predictive relationships among evidence, context,
and outcomes. These definitions presume neither an explicit causal graph nor a
general internal world model.

\subsection{Empirical Progression}
\label{sec:forecasting-as-inductive}

We ask three questions:
\begin{itemize}[nosep,leftmargin=*]
  \item \textbf{Do LMs react selectively to evidence?
  (\hyperref[sec:exp1]{Experiment~1}).}
  Across unresolved markets, every tested deployment distinguishes
  outcome-bearing packets from topical controls and moves in the withheld
  direction at high rates. A lexical baseline is directionally weaker than every
  deployment, establishing relevance-sensitive updating.

  \item \textbf{Can an LM select a relationship based on context?
  (\hyperref[sec:exp2]{Experiment~2}).}
  In Coin City, a controlled task, models forecast a target city's polls from
  response regimes displayed in two reference cities. Changing only the
  contextual cue changes which regime forecasts follow; direct target
  observations then weaken the cue.
  Episode-randomized labels test induction of a local convention, and open-model
  decoding and activation patching connect that distinction to behavior in the
  tested checkpoints.

  \item \textbf{Can finetuning transfer contextual relationships?
  (\hyperref[sec:exp3]{Experiments~3} and~\hyperref[sec:exp4]{4}).}
  Episode-matched training improves joint-shift transfer at two model scales.
  Separately, outcome-based training improves the tested bases on held-out
  historical markets. Scaling is
  nonmonotonic; several trained checkpoints reach the crowd-level floor.
\end{itemize}

Our contribution is a behavioral account of relationship use.
Across these settings, forecasts select, revise, and transfer context-dependent
relationships. Evaluation
should pair aggregate realism with contrasts that change the relationship while
holding easier cues fixed.
